\section{定义, 定理, 引理, 推论}

\begin{definition}[素数/质数]
    只能被 1 和自身整除且大于 1 的数
\end{definition}

\begin{definition}[合数]
    不是质数且大于 1 的数
\end{definition}

\begin{remark}
    1 既不是质数也不是合数
\end{remark}

\begin{theorem}
    \label{prime:infPrime}
    
    素数有无限多个
    
    \begin{proof}

        假设素数只有有限个,
        记作\(\def\enum#1{a_{ #1}}\enum{1},\enum{2},\dots,\enum{n}\)
        
        考虑\(N=\prod_{i=1}^np_i+1\), 显然\(N>2\)
        
        由于\(N\ne p_i,i=1,2,\dots,n\), 故必有\(N\)必有素因子\(p\)
        
        但\(p_i\nmid N,i=1,2,\dots,n\),
        故\(p\)是\(\def\enum#1{p_{ #1}}\enum{1},\enum{2},\dots,\enum{n}\)以外的素数,
        这与假设矛盾
        
    \end{proof}
\end{theorem}

\begin{lemma}[带余除法]
    \label{div:divWithRemaining}
    设\(a,b\in\mathbb{Z}\), 且\(b\geqslant1\), 则存在唯一的整数\(q,r\)使得
    \begin{equation}
        a=qb+r,r\in[0,b)\cap\mathbb{N}
    \end{equation}
    
    \begin{proof}
        首先, 取\(q=\lfloor\frac{a}{b}\rfloor\), \(r=a-qb\),
        容易验证此时的\(q,r\)满足条件
        
        下证唯一性, 假设又有一组整数\(q',r'\)使得\(a=q'b-r',~r'\in[0,b)\cap\mathbb{N}\),
        则
        
        \begin{equation}
            0=a-a=(q-q')b+(r-r')
        \end{equation}
        
        因此\(b\mid r-r'\)
        
        而\(|r-r'|\in[0,b)\cap\mathbb{N}\), 故只能有\(r-r'=0\), 即\(r=r'\)
        
        此时\(a-qb=a-q'b\), 有\(q=q'\)
        
    \end{proof}
    
\end{lemma}

\begin{theorem}
    \label{div:preBezout}
    设\(S\)为非空整数集, 若\(S\)关于整数的加减法封闭 \footnote{若\(a,b\in S\), 则\(a\pm b\in S\)}, 则存在唯一自然数\(d\)使得\(S=d\mathbb{Z}:=\{da|a\in\mathbb{Z}\}\) \footnote{条件中隐含了一点: 若\(a\in S\), 则\(ac\in S,~\forall c\in\mathbb{Z}\)} \footnote{一般的, 若 \(S\) 为主理想环, 则 \(\exists_1 d\in\mathbb{Z}~s.t.~S=\langle d\rangle\)}
    
    \begin{proof}

        若\(S=\{0\}\), 则只能取\(d=0\)
        
        设\(0\ne a\in S\), 则\(0=a-a\in S,-a=(-1)\times a\in S\),
        故此时\(S\)中必有正整数
        
        由良序公理可知\(S\)中的所有正整数中必有最小值,
        令\(d\)为\(S\)中的最小正整数, 下证\(S=d\mathbb{Z}\)
        
        首先易得\(d\mathbb{Z}\subseteq S\)
        
        之后我们在\(S\)中任取整数\(a\)做带余除法
        \begin{equation}
            a=qd+r,~q\in\mathbb{Z},r\in[0,d)\cap\mathbb{N}
        \end{equation}
        
        可知\(r=a-qd\in S\),
        而\(d\)为\(S\)中的最小正整数, 故必有\(r=0\)
        
        这表明\(a=qd\in d\mathbb{Z}\), 因此\(S\subseteq d\mathbb{Z}\)
        
        最后, 注意到\(d\)为\(S\)中的最小正整数,
        则满足\(S=d\mathbb{Z}\)的\(d\)一定是唯一的
        
    \end{proof}
\end{theorem}

\begin{corollary}[B\'{e}zout 定理]
    \label{gcdlcm:bezout}
    设\(\def\enum#1{a_{ #1}}\enum{1},\enum{2},\dots,\enum{n}\)是不全为零的整数, 且
    \((\def\enum#1{a_{ #1}}\enum{1},\enum{2},\dots,\enum{n})=d\), 则
    
    \begin{equation}
        S:=\left\{\sum_{i=1}^na_ix_i\bigg|x_i\in\mathbb{Z},i=1,2,\dots,n\right\}=d\mathbb{Z}
    \end{equation}
    
    \begin{proof}
        由定理 \ref{div:preBezout} 可知, 存在正整数\(a\)使得\(S=a\mathbb{Z}\), 接下来只需证\(a=d\)
        
        一方面, \(d\)是\(S\)中所有数的因子, 而\(a\in S\), 故\(d\mid a\)
        
        另一方面, 注意到\(\def\enum#1{a_{ #1}}\enum{1},\enum{2},\dots,\enum{n}\in S\), 故\(a\)是\(\def\enum#1{a_{ #1}}\enum{1},\enum{2},\dots,\enum{n}\)的公因子, 即有
        \begin{equation}
            a\mid(\def\enum#1{a_{ #1}}\enum{1},\enum{2},\dots,\enum{n})=d
        \end{equation}
        
        因此\(a=d\)
        
    \end{proof}
\end{corollary}
